%! Unit 3, Lessons 3.1–3.5 — Unit Test Study Guide (reference sheet; no answer key by design)
% This is a REFERENCE sheet, not a worksheet: every definition, every worked move, and every
% trap is already filled in, so there is nothing to answer and therefore no _key sibling.
% The practice a student does before the test is ../practice_test/, which HAS a key.
% NAMESTRIP: no \namedateperiod — the name row lives on a packet cover only.
% NOT GRADED: a study guide carries no \gradedblank and does not load grading (LESSON_SHAPE §9).
% Vocabulary below is quoted from the five lessons' own vocabulary boxes and notes rows
% (notes_key \termans entries) so a student reading this reads the same words twice.
% The full-credit sentences follow the test rubrics in ../../rubrics/.
% Worked contexts (camping festival shuttle; band warm-up; the car-sensor study) are deliberately
% NOT used in any Unit 3 lesson, AP practice, homework, or test.
%
% THE DATA (Move 8) 100 reshuffles, difference in mean seconds (drill - scales), step 0.5:
%   -2:1 -1.5:3 -1:10 -0.5:21 0:30 0.5:21 1:11 1.5:2 2:1  (sum 100).
%   Observed 1.5: 2+1 = 3 of 100 at or above -> significant.
%   Observed 0.5: 21+11+2+1 = 35 of 100 at or above -> not significant.
\documentclass[12pt]{article}
\usepackage{apstats-article}
\usepackage{apstats-boxes}
\usepackage{ltablex}
\keepXColumns
% One tabularx per group of lessons rather than one long table: a tabularx with X
% columns is an unbreakable box, so a single long table floats whole to a fresh page
% and strands most of page 1. Short tables flow, and each carries its lesson tag.

% Optional arg: the \boxguard size (lines kept with the strip).
\newcommand{\sghead}[2][10]{%
  \vspace{0.13in}
  \boxguard[#1]
  \begin{headlinebox}{navy}\color{white}\bfseries\small #2\end{headlinebox}
  \vspace{0.09in}}

\begin{document}

\pageheader{Unit 3 \quad Lessons 3.1--3.5}{Unit Test Study Guide}

\vspace{0.12in}

\boxguard[8]
\begin{remindbox}
\small
The test covers \textbf{Lesson 3.1} through \textbf{Lesson 3.5} --- how data are collected and
what they let you conclude. Ten multiple-choice items (two per lesson), then two free-response
questions: one on \textbf{a survey sample}, one on \textbf{an experiment} with a reshuffle
dotplot already drawn. 20 points. Every answer must be \textbf{in context} --- name the people or
things studied and the variable measured.
\end{remindbox}

\sghead{1. Vocabulary you are expected to use, not just recognize}

\small
\renewcommand{\arraystretch}{1.15}

\textbf{From Lesson 3.1 --- planning a study}
\par\vspace{3pt}
\noindent\begin{tabularx}{\linewidth}{@{} l X @{}}
\rowcolor{sky}
\textbf{Term} & \textbf{What it means} \\
\midrule
Population & the entire group of individuals we want to know about. \\
Sample & the part (subset) of the population we actually measure. \\
Census & data from \textbf{every} individual in the population. \\
Observational study & observe and measure, without imposing a treatment. \\
Experiment & researchers assign treatments and measure the response. \\
\end{tabularx}

\vspace{7pt}
\boxguard[9]
\textbf{From Lesson 3.2 --- sampling methods}
\par\vspace{3pt}
\noindent\begin{tabularx}{\linewidth}{@{} l X @{}}
\rowcolor{sky}
\textbf{Term} & \textbf{What it means} \\
\midrule
Simple random sample & every \textbf{group} of $n$ individuals has an equal chance to be the sample. \\
Without replacement & an individual can be chosen only once. \\
Stratified random sample & an SRS from \textbf{every} stratum (groups alike inside), combined. \\
Cluster sample & randomly choose \textbf{some} clusters (groups that each look like the
                   whole); measure \textbf{everyone} in them. \\
\end{tabularx}

\vspace{7pt}
\boxguard[9]
\textbf{From Lesson 3.3 --- what goes wrong}
\par\vspace{3pt}
\noindent\begin{tabularx}{\linewidth}{@{} l X @{}}
\rowcolor{sky}
\textbf{Term} & \textbf{What it means} \\
\midrule
Convenience sample & whoever is easiest to reach. No chance is involved. \\
Voluntary response bias & people choose to respond; those with strong opinions are
                   over-represented. \\
Undercoverage bias & part of the population has a reduced (or no) chance of being chosen. \\
Nonresponse bias & chosen individuals can't be reached or refuse, and differ from those who
                   answer. \\
Response bias & the question, or how it is asked, pushes the answers one way. \\
\end{tabularx}

\vspace{7pt}
\boxguard[6]
\textbf{From Lesson 3.4 --- the parts of an experiment}
\par\vspace{3pt}
\noindent\begin{tabularx}{\linewidth}{@{} l X @{}}
\rowcolor{sky}
\textbf{Term} & \textbf{What it means} \\
\midrule
Experimental units & the individuals that are assigned the treatments (\emph{subjects} if people). \\
Treatment & a level (or combination of levels) of a factor, set on purpose. \\
Response variable & what is measured \emph{after} the treatments. \\
\end{tabularx}

\vspace{7pt}
\boxguard[12]
\textbf{Also from Lesson 3.4 --- what makes the design work}
\par\vspace{3pt}
\noindent\begin{tabularx}{\linewidth}{@{} l X @{}}
\rowcolor{sky}
\textbf{Term} & \textbf{What it means} \\
\midrule
Confounding variable & a variable tied to the explanatory variable that also affects the
                   response, so the two effects get mixed up. \\
Control group & units given no treatment or a placebo, as a baseline to compare against. \\
Placebo & an inactive look-alike treatment; the \emph{placebo effect} is a response to it. \\
Single- / double-blind & single: one side does not know who gets which treatment; double:
                   neither subjects nor staff know. \\
Four principles & \textbf{comparison}, \textbf{random assignment}, \textbf{replication},
                   \textbf{control}. \\
Randomized block & sort units into blocks alike on something that affects the response, then
                   randomly assign \emph{within} each block. \\
Matched pairs & a block design whose blocks have two units: chance picks who gets which, or
                   each unit gets both in a random order. \\
\end{tabularx}

\vspace{7pt}
\boxguard[9]
\textbf{From Lesson 3.5 --- inference and experiments}
\par\vspace{3pt}
\noindent\begin{tabularx}{\linewidth}{@{} l X @{}}
\rowcolor{sky}
\textbf{Term} & \textbf{What it means} \\
\midrule
Statistical inference & drawing conclusions beyond the individuals we measured. \\
Statistically significant & a result too large to be explained by chance alone. \\
Scope of inference & what a study lets us conclude: cause and effect, and about whom. \\
Generalize & apply the results to the population the individuals were picked from. \\
\end{tabularx}

\normalsize

\sghead{2. The moves the test will ask you to make}

\small
\textbf{Move 1 --- Observational study or experiment? (3.1)} Ask one question: \emph{did the
researchers impose a treatment?} \emph{Worked:} an insurer puts a sensor in 2{,}000 cars and
logs every hard stop for a year; drivers who use their phones more have more hard stops. The
measuring is precise, but nobody was assigned phone use, so it is \textbf{observational}.

\vspace{3pt}
\textbf{Move 2 --- A sample speaks only for its population (3.1).} \emph{Worked:} if those 2{,}000
drivers were picked at random from the insurer's \textbf{Texas} customers, the result speaks for
the insurer's Texas customers --- not all drivers, and not just the 2{,}000.

\vspace{3pt}
\textbf{Move 3 --- Choose an SRS, and test whether a plan is one (3.2).} A camping festival has
12{,}000 attendees, in 300 rows of 40 tents, and wants the proportion who would pay for a shuttle
to town. \emph{Worked:} \textbf{label} the attendees 1 to 12{,}000; \textbf{randomize} with a
random number generator; \textbf{select} 200 different labels, skipping repeats. Picking 5 rows at
random instead gives everyone the same $\tfrac{5}{300}$ chance, but it is \textbf{not} an SRS: a
group of 200 spread across more than 5 rows can never be the sample.

\vspace{3pt}
\boxguard[9]
\textbf{Move 4 --- Stratified or cluster? (3.2)} Are the groups alike inside, or does each one
look like the whole?

\vspace{2pt}
\noindent\begin{tabularx}{\linewidth}{@{} l X X @{}}
\rowcolor{sky}
 & \textbf{Stratified} & \textbf{Cluster} \\
\midrule
The groups & alike inside (age groups) & each a mix (rows of tents) \\
What you do & an SRS from \textbf{every} group & \textbf{all} of a few random groups \\
Why & more precise when the variable differs across groups & cheaper and easier to reach \\
\end{tabularx}

\vspace{5pt}
\boxguard[6]
\textbf{Move 5 --- Name the bias, then give its direction (3.3).} Three parts, every time: the
\textbf{name}; \textbf{who} is over- or under-represented, in context; \textbf{too high or too
low}. \emph{Worked,} all estimating the proportion who would pay for the shuttle:

\vspace{2pt}
\noindent\begin{tabularx}{\linewidth}{@{} l X l @{}}
\rowcolor{sky}
\textbf{Bias} & \textbf{What happened, and who it favors} & \textbf{Estimate} \\
\midrule
Voluntary response & An online poll: 4{,}000 vote. Those fed up with parking vote most. & too high \\
Undercoverage & The list holds only app buyers. Older attendees bought at the gate, and they
                want a ride more. & too low \\
Nonresponse & Of an SRS of 200, 90 reply. Those without cars have the most reason to. & too high \\
Response & ``Wouldn't a shuttle end the dangerous parking chaos?'' leads people to yes. & too high \\
\end{tabularx}

\vspace{3pt}
None of these is cured by size: \textbf{a larger sample reduces variability, not bias.} The
200-person SRS beats the 4{,}000 votes.

\vspace{3pt}
\textbf{Move 6 --- Lay out an experiment (3.4).} A director tests whether a 10-minute breathing
drill lengthens how long band students can hold a note. 48 volunteers are randomly assigned, 24
to the drill and 24 to a 10-minute scales warm-up. \emph{Worked:} \textbf{units:} the 48
students. \textbf{Treatments:} drill, scales. \textbf{Response:} seconds the note is held.
Random assignment balances variables like practice habits, so a difference can be traced to the
drill. If the students know what they did but the person timing does not, it is
\textbf{single-blind}.

\vspace{3pt}
\textbf{Move 7 --- Block or pair, and keep the chance (3.4).} Block only on what affects the
response. \emph{Worked:} brass and woodwind players hold notes differently, so form two blocks and
\textbf{randomly assign within each}. Birth month would be a useless block. In a matched pairs
version, each student does both warm-ups on different days and \textbf{a coin flip decides which
comes first} --- otherwise a practice effect always favors one treatment.

\vspace{3pt}
\boxguard[14]
\textbf{Move 8 --- Read the reshuffle dotplot (3.5).} Each dot is one reshuffle: one difference
that chance alone produced, as if the treatment did nothing.

\vspace{2pt}
\noindent\begin{minipage}[c]{0.43\linewidth}\raggedright
\emph{Worked:} the drill group's mean was \textbf{1.5 seconds longer}. Count the dots at 1.5 or
more: $2+1=$ \textbf{3 of 100}. Chance rarely does that: \textbf{statistically significant}. Had
it been only 0.5 seconds, $21+11+2+1=$ \textbf{35 of 100} reach it: chance could easily explain
it, so \textbf{not} significant.
\end{minipage}\hfill
\begin{minipage}[c]{0.54\linewidth}\centering
\begin{tikzpicture}[x=0.9cm, y=0.105cm]
  \draw[thick] (-4.5,0) -- (4.5,0);
  \foreach \x/\l in {-4/-2.0,-3/-1.5,-2/-1.0,-1/-0.5,0/0,1/0.5,2/1.0,3/1.5,4/2.0} {
    \draw (\x,0) -- (\x,-1.6); \node[below, font=\tiny] at (\x,-1.4) {$\l$};}
  \foreach \x/\n in {-4/1,-3/3,-2/10,-1/21,0/30,1/21,2/11,3/2,4/1}
    \foreach \k in {1,...,\n} \fill[navy] (\x,\k) circle (1.05pt);
  \draw[redacc, dashed, thick] (3,3.5) -- (3,13);
  \node[font=\tiny\bfseries, redacc, above] at (3,13) {observed 1.5};
  \node[font=\tiny] at (0,-7.5) {Simulated difference in mean seconds, drill $-$ scales};
  \node[font=\scriptsize\bfseries] at (0,34) {100 random reshuffles};
\end{tikzpicture}
\end{minipage}

\vspace{5pt}
\boxguard[12]
\textbf{Move 9 --- Scope of inference: two separate questions (3.5).} Was the treatment
\emph{assigned} at random? That gives \textbf{cause}. Were the individuals \emph{selected} at
random? That lets you \textbf{generalize}.

\vspace{3pt}
{\centering\footnotesize
\renewcommand{\arraystretch}{1.25}
\begin{tabular}{|c|c|c|}
\hline
 & \cellcolor{sky}\textbf{Selected at random} & \cellcolor{sky}\textbf{Not selected at random} \\
\hline
\cellcolor{sky}\textbf{Assigned at random} & Cause: yes \quad Generalize: yes & Cause: yes \quad Generalize: no \\
\hline
\cellcolor{sky}\textbf{Not assigned} & Cause: no \quad Generalize: yes & Cause: no \quad Generalize: no \\
\hline
\end{tabular}\par}
\normalsize

\vspace{0.10in}
\boxguard[20]
\begin{scenariobox}[The sentences that earn full credit]{navy}
\small
\textbf{Choose an SRS:} ``Label the attendees 1 to 12{,}000. Generate random integers from 1 to
12{,}000, ignoring repeats, until 200 different labels are chosen. Survey those 200
attendees.''\par\vspace{3pt}
\textbf{Why stratified is more precise:} ``Interest in a shuttle is likely related to age, so each
age group is alike inside and every age group is represented, giving a less variable
estimate.''\par\vspace{3pt}
\textbf{Bias and direction:} ``Nonresponse bias. Attendees without cars have the most reason to
reply, so they are over-represented, and the estimated proportion who would pay is likely too
high.''\par\vspace{3pt}
\textbf{Random assignment with slips:} ``Write each student's name on one of 48 identical slips,
mix them in a hat, and draw 24 without replacement. Those 24 do the drill; the other 24 do
scales.''\par\vspace{3pt}
\textbf{Blocking:} ``Separate the students into a brass block and a woodwind block. Within each
block, randomly assign half to the drill and half to scales, and compare within blocks.''
\par\vspace{3pt}
\textbf{Significance:} ``Only 3 of the 100 reshuffles gave a difference of 1.5 seconds or more, so
a difference that large rarely happens from the random assignment alone. It is statistically
significant.''\par\vspace{3pt}
\textbf{Cause and to whom:} ``Yes, the drill caused the longer notes, because the treatments were
randomly assigned and the difference is significant. But the students were volunteers, not
randomly selected, so the result applies only to students like these.''\par\vspace{3pt}
Each names the people, the treatment, and the variable. Without context, it is not an answer.
\end{scenariobox}

\sghead{3. Eight traps that cost points}

\small
\begin{enumerate}[leftmargin=1.6em, itemsep=2pt]
  \item \textbf{Careful measuring is not a treatment.} Trackers, apps, and sensors make a study
        precise; only \emph{imposing} a treatment makes it an experiment.
  \item \textbf{Equal chance for each person is not an SRS.} Every \emph{group} of $n$ must have
        an equal chance.
  \item \textbf{Stratified: some from every group. Cluster: all from some groups.} Groups alone
        settle nothing.
  \item \textbf{Random samples can still be biased.} A random sample from a bad list, with people
        who never answer, or with a leading question is still wrong.
  \item \textbf{A bigger sample does not fix bias.} Size reduces variability only. Tens of
        thousands of volunteers are still volunteers.
  \item \textbf{Blocking and pairing do not replace chance.} Randomize inside every block, and
        for every pair. Never let a fixed rule or the subject choose.
  \item \textbf{Not every difference is significant.} Count the dots at or beyond the observed
        value, on its side of the plot. Many dots means chance can explain it.
  \item \textbf{A random sample is not a random assignment.} Selection lets you generalize;
        only assignment lets you say cause. Volunteers can show cause, but only for those like
        them.
\end{enumerate}
\normalsize

\sghead[6]{4. Before you walk in}

\small
Work the \textbf{sample test} without notes, then check it against its key. A free-response
answer shorter than the key's is usually missing the context. Re-read each \textbf{Keep in Mind} box.
\normalsize

\end{document}
